\documentclass[thesis=da, language=german, paper=a4]{isys}[2013/06/28]
%
% [2013/06/28] ensure that the newest isys.cls is used
%              if an older version is used, latex issues a warning
%
% thesis   = sa       - Studienarbeit  / Student thesis 
%            da       - Diplomarbeit   / Diploma thesis
%            bachelor - Bachelorarbeit / Bachelor thesis
%            master   - Masterarbeit   / Master thesis
%            semester - Semesterarbeit / Semester thesis
%
% language = {english, german} 
%            Set default language of document. The class file isys.cls loads the babel
%              package with both options [english, ngerman]. For text that shall be 
%              typeset in a different language than the default language, the 
%              environment 'otherlanguage' or the command '\foreignlanguage' may be used
%              with either of these options (see babel package documentation)
%
% paper    = {a4, letter} - paper size
%
% thesisdraft - mark document as draft without passing "draft" option to 
%                 other packages (e.g. pictures are shown in the document, etc.)
%
% oneside - this option avoids empty pages in the document output 
%                which are necessary for twoside printing, which is useful for 
%                draft versions. 
%                DO NOT USE THIS OPTION WHEN SUBMITTING YOUR THESIS!
%

% mark overful boxes by a 2pt wide rule 
% -- FIX ALL OVERFULL BOXES BEFORE SUBMITTING YOUR THESIS!
\overfullrule=2pt 


%%%%%%%%%%%%%%%%%%%%%%%%%%%%%%%%%%%%%%%%%%%%%%%%%%%%%%%%%%%%%%%%%%%
% Load packages - useful packages, adapt to your own needs
%
% isys.cls loads the "graphicx" package, so \includegraphics can 
%  directly be used to include figures.
%
%%%%%%%%%%%%%%%%%%%%%%%%%%%%%%%%%%%%%%%%%%%%%%%%%%%%%%%%%%%%%%%%%%%
% input encoding - UTF8 or latin1 allow directly typing german umlauts
% make sure the file encoding of all files matches this encoding!
\usepackage[utf8]{inputenc}
% font enconding - do not change this unless you exactly know what you're doing!
\usepackage[T1]{fontenc}

% automatically use correct quotes for the current language: \enquote{text in quotes}
\usepackage{csquotes}
\MakeAutoQuote{»}{«}  % input »text in quotes« is equivalent to \enquote{text in quotes}
\let\q=\enquote 			% use \q{text in quotes}, which is shorter than \enquote{text in quotes}

% better math-commands (e.q.: \dddots, \eqref{})
\usepackage{amsmath,amssymb}                                 

% better-looking rules in tabulars, \toprule, \midrule, \bottomrule
\usepackage{booktabs}

% use consistent style for units (upright shape)
% \unit[value]{unit}, \unitfrac[value]{unit_numerator}{unit_denominator}
\usepackage{units}

% tables with fixed table width 
\usepackage{tabularx, tabulary}
% improve ragged right alignment with \RaggedRight
\usepackage{ragged2e}
% enable import of eps pictures using pdflatex
%\usepackage{epstopdf}

% enable hyperlinks and bookmarks in PDF file (must typically be loaded as LAST package)
% if used, print hyperlinks in black and suppress boxes around them 
\usepackage[colorlinks=true, allcolors=black, bookmarksnumbered=true]{hyperref}
% to set bookmarks on the first latex run (hyperref needs two),
% provides addtional options for bookmarks (use carefully!)
\usepackage{bookmark}


% Bibliography using classical bibtex
%\bibliographystyle{alphadin}

% better: use biblatex! provides much more options
% use \printbibliography to print the bibliography 
% (search for \printbibliography in this document and uncomment the line, 
%  comment out the line for bibtex)
\usepackage[style=alphabetic,
						backend=bibtex8,		% or use biber or bibtex?
						bibencoding=latin1,	% .bib file in latin1 or utf8 encoding?
						maxbibnames=99			% give full list of authors (up to 99) in bibliography
						]{biblatex}
\bibliography{references}


%%%%%%%%%%%%%%%%%%%%%%%%%%%%%%%%%%%%%%%%%%%%%%%%%%%%%%%%%%%%%%%%%%%
% Information for title page 
%%%%%%%%%%%%%%%%%%%%%%%%%%%%%%%%%%%%%%%%%%%%%%%%%%%%%%%%%%%%%%%%%%%

% title in the language in which the thesis is written
\title{Schätzverfahren für adaptive Regelung von Hybridgetrieben}
%\subtitle{ }
% title in foreign language (german or english, depending on the language of the thesis)
\foreigntitle{Estimation Procedure for Adaptive Control of Hybrid Gearboxes}
%\foreignsubtitle{Estimation Procedure for Adaptive Control of Hybrid Gearboxes}

\author{Koray Karci}
\major{Maschinenbau}
\date{30.09.2018}		% use either ISO format YYYY-MM-DD or german format DD.MM.YYYY
										% default: current date

\examiner{Prof. Dr.-Ing. Oliver Sawodny}
%\examiner{Prof. Dr.-Ing. Cristina Tar\'{i}n}
%\supervisor{Dipl.-Ing. Supervisor 1}		% give max. 4 \supervisor {...} commands
\supervisor{M.Sc. Michel Bauer}
%\supervisor{Dipl.-Ing. Supervisor 3}
%\supervisor{M.Sc. Supervisor 4}

%% if work on thesis has been conducted with a partner company/institute
%% use \partner and (if desired also \partnerLogo)
% \partnerLogo has no effect as long as no \partner{...} is given
%\partner{Example company name\\
%			Second line company name}
%\partnerLogo{\includegraphics[height=1.4cm, clip, viewport=0 0 30 10]{logos/logoISYS_deutsch}}



%%%%%%%%%%%%%%%%%%%%%%%%%%%%%%%%%%%%%%%%%%%%%%%%%%%%%%%%%%%%%%%%%%%
% Start of document output
%%%%%%%%%%%%%%%%%%%%%%%%%%%%%%%%%%%%%%%%%%%%%%%%%%%%%%%%%%%%%%%%%%%
\begin{document}

% Titelblatt ausgeben / print title page
\pdfbookmark{Titelseite}{title}
\maketitle

%% Englisches Titelblatt in deutschsprachiger Arbeit erstellen
%\begin{otherlanguage}{english}\maketitle \end{otherlanguage}

%% Create german title page in english theses
%\begin{otherlanguage}{ngerman}\maketitle \end{otherlanguage}


	
% Kurzfassung in der Dokumentensprache / Abstract in default document language
\pdfbookmark{Kurzfassung}{abstract-de}
\begin{abstract}
Kurzfassung
\end{abstract}

% Falls gewuenscht, englischen Abstract in deutschsprachiger Arbeit einbinden
\newpage
\begin{otherlanguage}{english}
\pdfbookmark{Abstract}{abstract-en}
\begin{abstract}
Abstract 
\end{abstract}
\end{otherlanguage}

%% Inhalts-/Abbildungs-/Tabellenverzeichnis
%% Table of contents, List of Figures/Tables

	
	
%% If subsections shall also be included in the table of contents
%% uncomment next line. Default: only sections are included.
\setcounter{tocdepth}{1} % include only chapters and sections (1.1) - default
%\setcounter{tocdepth}{2} % also include subsections - use only if told to do so by your advisor!
\cleardoublepage
\pdfbookmark{\contentsname}{toc}
\tableofcontents

\addchap{Abkürzungen und Formelzeichen}

\section*{Abkürzungen}

\begin{flushleft}
\begin{tabularx}{\textwidth}{XlX}	
	&&\\
	DOF								&Freiheitsgrade (Engl.: Degrees of Freedom)&\\
	E/A								&Eingang/Ausgang &\\
	IMC								&Internal Model Control&\\
	MPC								&Modellprädiktive Regelung (Engl.: Model Predictive Control)&
\end{tabularx}
\end{flushleft}


\vspace{20mm}
\section*{Formelzeichen}


\begin{flushleft}
\begin{tabularx}{\textwidth}{llX}
	\toprule
	Formelzeichen					 	& Einheit									& Beschreibung \\
	\midrule
	$I$											&	A												&	Strom \\
	$U$											&	V												&	Spannung \\
	$T$											&	K												&	Temperatur \\
	&&\\
	$\tau$									&	s												& Zeitkonstante \\
	$\rho_{\mathrm{L}}$			&	$\unitfrac[]{kg}{m^3}$	& Dichte der Luft \\
	&&\\
	$\boldsymbol{u}$						& 												&	Regeleingang des zu regelnden Systems \\
	$\boldsymbol{x}$						& 												&	Zustandsvektor \\
	$\boldsymbol{y}$						& 												&	Regelausgang des zu regelnden Systems \\
\bottomrule
\end{tabularx}
\end{flushleft}


%% Chapters and sections without numbering but included in running heading
%\addchap{Chapter without numbering, included in table of contents}
%\addchap*{Chapter without numbering, not included in table of contents}
%\addsec{Section without numbering, included in table of contents}
%\addsec*{Section without numbering, not included in table of contents}

%%%%%%%%%%%%%%%%%%%%%%%%%%%%%%%%%%%%%%%%%%%%%%%%%%%%%%%%%%%%%%%%%%%
% Beginn des Textes Inhalts / Start of thesis content
%%%%%%%%%%%%%%%%%%%%%%%%%%%%%%%%%%%%%%%%%%%%%%%%%%%%%%%%%%%%%%%%%%%

	
% It may be helpful to put each chapter into a separate file. 
% \includeonly{...} allows to restrict the files that are included,
% which may speed up the document creation process.
% \include always starts a new odd page before including the file 
% contents. If this is not desired, \input can be used.
\chapter{Einleitung und Motivation}\label{ch:ch1}
Dies ist die Einleitung und Motivation.




\chapter{Grundlagen, Theorie, Vorarbeiten}\label{ch:ch2}
Die Grundlagen werden hier erläutert. Das ist eine Änderung



\chapter{Hauptteil}\label{ch:ch3}
Der Hauptteil kann aus mehreren Kapiteln bestehen. Hier soll auf die Vorgehensweise sowie die Ergebnisse der Arbeit eingegangen werden. Dabei ist auf eine schlüssige und klare Strukturierung zu achten, die den roten Faden der Arbeit verdeutlicht. 

Kurzer einleitender Text nach jedem Kapitelbeginn und vor Beginn der \texttt{sections}. Es stehen also nie zwei Überschriften direkt hintereinander. An dieser Stelle soll dem Leser ein Überblick über das folgende Kapitel gegeben werden.

Hier habe ich etwas geändert
\section{Erstes Unterkapitel}\label{sec:ch3_foobar}
Dies wurde bereits von Fliess et~al. im Jahre 1995 gezeigt~\cite{Fliess.1995}.

Mit \texttt{biblatex}-Befehlen: Dies wurde bereits von \citeauthor{Fliess.1995} im Jahre~\citeyear{Fliess.1995} gezeigt~\cite{Fliess.1995}.


Das Logo des Instituts ist in Abbildung~\ref{fig:logo_isys1} dargestellt. 

  
\begin{figure} % avoid using additional specifiers for graphics position except for final tuning
 \centering    % preferable over \begin{center}...\end{center} as the latter includes space before and after the figure 
 \includegraphics[width=0.5\textwidth]{logos/logoISYS_deutsch}
 \caption{Put figure captions \emph{below} the figure.}
 \label{fig:logo_isys1} % always use \label AFTER \caption
\end{figure}

Die erzielten Ergebnisse sind in Tabelle~\ref{fig:example_table1} zusammengefasst.

\begin{table} % avoid using additional specifiers for graphics position except for final tuning
	\caption{Put table captions \emph{above} the table. Put units in header. Avoid vertical lines in a table.}
	\label{fig:example_table1}
	\centering
	\begin{tabular}{lll}
		\toprule
		Header & Header (\unit{unit}) & Header (\unitfrac{unit}{frac})   \\
		\midrule
		Content 1 & Content 2 & Content 3 \\
		Content 1 & Content 2 & Content 3 \\
		\bottomrule
	\end{tabular}
\end{table}

   

\section{Zweites Unterkapitel}


\subsection{Neuer Abschnitt}





\chapter{Zusammenfassung und Ausblick}\label{ch:ch4}
Zusammenfassung und Ausblick werden hier beschrieben.
\cleardoublepage

\appendix
\chapter{Anhang}\label{ch:appCh1}
Im Anhang werden Ergebnisse aufgeführt, die zwar im Kontext der Arbeit wesentlich, jedoch nicht für das Verständnis der Arbeit notwendig sind. Beispielsweise sind dies detailliertere Messergebnisse, Herleitungen oder Datenblätter. 

% List of tables and list of figures at end
\listoffigures
\listoftables

% Literaturverzeichnis
% classical bibtex
%\bibliography{references}

% modern biblatex
\printbibliography[heading=bibintoc]

% in any case ...
\label{ch:bib}



% Erklaerung des Autors, netig fuer Bachelor-/Master-Arbeiten der Universitaet Stuttgart
% Declaration of authorship, necessary for bachelor and master thesis at University of Stuttgart
\DeclarationOfAuthorship
\end{document}
